\section{Kết quả, bài học và hướng phát triển}

\subsection{Mức độ đáp ứng mục tiêu}

Đối chiếu mục tiêu ban đầu với MVP và evidence. Tách rõ mục tiêu đạt, đạt một phần và chưa đạt; giải thích nguyên nhân kỹ thuật hoặc quản lý.

\subsection{Bài học của nhóm}

Tập trung vào việc quản lý requirement, đồng bộ artefact, làm việc nhóm, tích hợp code, kiểm thử, xử lý dữ liệu IoT, demonstration và sử dụng AI có trách nhiệm. Nêu tình huống cụ thể thay vì kết luận chung chung.

\subsection{Giới hạn và hướng phát triển}

Phân biệt giới hạn do thời gian/phạm vi môn học với hạn chế kiến trúc cần cải thiện. Chỉ đề xuất feature tương lai khi đã chỉ ra nhu cầu và tác động đến requirement, design, implementation và test.

\section{Phụ lục}

\subsection{Tóm tắt biên bản họp}
Liệt kê buổi họp, ngày giờ, quyết định và đường dẫn tới biên bản của từng phase. Không chèn toàn bộ biên bản nếu làm báo cáo mất trọng tâm, trừ khi LMS yêu cầu hợp nhất.

\subsection{Hướng dẫn truy cập source và demonstration}
Ghi repository/commit/tag dùng để chấm, tài khoản demo, URL hoặc gói Docker, video dự phòng và checksum/version nếu cần.

\subsection{Bảng truy vết cuối}
Đối chiếu tuần tự stakeholder need và user story với requirement, use case, screen, behavioral diagram, component, class, method, test và demo step. Bảng này dùng để phát hiện artefact mồ côi trước khi nộp.
